\ZSFNewColumn
\StartChapter{Lineare Abbildungen}

\begin{runintext}
Eine \ZSFkeyword*{lineare Abbildung} $\mathcal{F} \colon V \to W$ ist eine Abbildung zwischen zwei Vektorräumen über $\R$, die Addition und Skalarmultiplikation erhält.
\end{runintext}

\SubsectionBar{Definition}
\ZSFindex{Lineare Abbildung}

\begin{defbox}[Lineare Abbildung]
Eine Abbildung $\mathcal{F}$ heisst linear, wenn für alle Vektoren $x, y \in V$ und jeden Skalar $\alpha \in \R$ gilt:
  \begin{ZSFlist}[][ordered]
    \ZSFItem{Additivität: $\mathcal{F}(x + y) = \mathcal{F}(x) + \mathcal{F}(y)$.}
    \ZSFItem{Homogenität: $\mathcal{F}(\alpha \cdot x) = \alpha \cdot \mathcal{F}(x)$.}
  \end{ZSFlist}
\end{defbox}

\begin{runintext}
\ZSFconclusion{Folgerungen:}
  \begin{ZSFlist}
    \ZSFItem{Jede lineare Abbildung bildet den Nullvektor \ZSFdanger{immer} auf den Nullvektor ab: $\mathcal{F}(0_V) = 0_W$.}
    \ZSFItem{Abbildungen zwischen endlichdimensionalen Vektorräumen lassen sich stets durch Matrixmultiplikation $x \mapsto A \cdot x$ darstellen.}
  \end{ZSFlist}
\end{runintext}

\SubsectionBar[sec:kernel-image]{Kern und Bild \ZSFref{sec:hub-kernel-image-dimension}}

\ZSFindexsee{Nullraum}{Kern}
\ZSFindexsee{Spaltenraum}{Bild}
\ZSFindexsee{Rangsatz}{Dimensionssatz}
\ZSFindexsee{Ker}{Kern}
\ZSFindexsee{Im}{Bild}
\begin{defbox}[Kern (Lösungsraum)]
Der \ZSFkeyword{Kern} $\Ker(A)$ von $A \in \R^{m \times n}$ enthält alle Vektoren $x$, die auf den Nullvektor abgebildet werden:
\[ \Ker(A) = \{x \in \R^n \mid A \cdot x = 0 \} \]
Die Dimension des Kerns ist gleich der Anzahl der freien Parameter im homogenen LGS:
\[ \dim(\Ker(A)) = n - \rang(A) \]
Um den Kern zu bestimmen, löst man das LGS $A \cdot x = 0$.
\end{defbox}

\begin{defbox}[Bild (Spaltenraum)]
Das \ZSFkeyword{Bild} $\Im(A)$ von $A \in \R^{m \times n}$ ist die Menge aller möglichen Funktionswerte:
\[ \Im(A) = \{y \in \R^m \mid \exists x \in \R^n \colon y = A \cdot x \} \]
Das Bild wird von den Spaltenvektoren von $A$ aufgespannt: $\Im(A) = \spanop(a_1, \dots, a_n)$. \ZSFref{sec:basis}
\end{defbox}

\begin{defbox}[Zusammenhänge und Dimensionssatz][weight=quiet]
Für jede Matrix $A \in \R^{m \times n}$ gilt:
  \begin{ZSFlist}
    \ZSFItem{\ZSFkeyword{Dimensionssatz} (Rangsatz):
  \[ \dim(\Im(A)) + \dim(\Ker(A)) = n \]}
    \ZSFItem{Spaltenrang ist gleich Zeilenrang:\ZSFbreak $\dim(\Im(A)) = \rang(A)$.}
    \ZSFItem{\ZSFkeyword{Orthogonales Komplement}: $\Im(A) \perp \Ker(A^T)$. \ZSFref{sec:orthogonality-projection}}
    \ZSFItem{\ZSFkeyword{Fredholm-Alternative}: $Ax = b$ ist lösbar $\Longleftrightarrow b \in \Im(A) \Longleftrightarrow b \perp \Ker(A^T)$.}
  \end{ZSFlist}
\end{defbox}

\ZSFNewColumn
\SubsectionBar[sec:map-matrix]{Bestimmung der Abbildungsmatrix}
\ZSFindex{Abbildungsmatrix}
\ZSFindexsee{Darstellungsmatrix}{Abbildungsmatrix}

\begin{runintext}
Um eine Matrix $A$ für eine lineare Abbildung $\mathcal{F} \colon V \to W$ zu bestimmen, bilden wir die Basisvektoren von $V$ ab und stellen die Ergebnisse als Spalten von $A$ dar.
\newline
Beispiel: $\mathcal{F} \colon P_2 \to P_1$, $p(x) \mapsto p'(x)$ (Ableitung).
\end{runintext}

\begin{ZSFlist}[Konstruktion der Abbildungsmatrix][ordered]
  \ZSFItem{Basen festlegen}{Wähle Basen für Startraum $V$ und Zielraum $W$ \ZSFref{sec:basis}.
  \newline
  Startbasis $P_2$: $\mathcal{B}_V = \{x^2, x, 1\}$.
  \newline
  Zielbasis $P_1$: $\mathcal{B}_W = \{x, 1\}$.}
  \ZSFItem{Basisvektoren abbilden}{Berechne die Bilder der Basisvektoren und stelle sie als Koordinaten bezüglich der Zielbasis $\mathcal{B}_W$ dar \ZSFref{sec:coordinates}:
  \newline
  $\mathcal{F}(x^2) = (x^2)' = 2x = 2 \cdot x + 0 \cdot 1 \implies \begin{pmatrix} 2 \\ 0 \end{pmatrix}$
  \newline
  $\mathcal{F}(x) = (x)' = 1 = 0 \cdot x + 1 \cdot 1 \implies \begin{pmatrix} 0 \\ 1 \end{pmatrix}$
  \newline
  $\mathcal{F}(1) = (1)' = 0 = 0 \cdot x + 0 \cdot 1 \implies \begin{pmatrix} 0 \\ 0 \end{pmatrix}$}
  \ZSFItem{Matrix aufstellen}{Trage die Koordinatenvektoren als Spalten in die Matrix $A$ ein:
  \[ A = \begin{pmatrix} 2 & 0 & 0 \\ 0 & 1 & 0 \end{pmatrix} \]}
\end{ZSFlist}

\SubsectionBar[sec:rotation-matrices]{Drehmatrizen}
\ZSFindex{Drehmatrix}
\ZSFindexsee{Rotationsmatrix}{Drehmatrix}

\begin{runintext}
Drehungen im $\R^3$ werden durch orthogonale Matrizen dargestellt \ZSFref{sec:orthogonal-matrices}. Die Drehrichtung folgt der \ZSFkeyword{Rechte-Hand-Regel}:
\end{runintext}

\begin{formulabox}[Drehung um die x-Achse]
  \[
  R_x(\varphi) = \begin{pmatrix}
    1 & 0 & 0 \\
    0 & \cos\varphi & -\sin\varphi \\
    0 & \sin\varphi & \cos\varphi
  \end{pmatrix}
  \]
  \ZSFsep[Drehung um die y-Achse]
  \[
  R_y(\varphi) = \begin{pmatrix}
    \cos\varphi & 0 & \sin\varphi \\
    0 & 1 & 0 \\
    -\sin\varphi & 0 & \cos\varphi
  \end{pmatrix}
  \]
  \ZSFsep[Drehung um die z-Achse]
  \[
  R_z(\varphi) = \begin{pmatrix}
    \cos\varphi & -\sin\varphi & 0 \\
    \sin\varphi & \cos\varphi & 0 \\
    0 & 0 & 1
  \end{pmatrix}
  \]
\end{formulabox}
